%% GENERATED FILE -- do not hand-edit.
%% SOURCE: experiments/database/scripts/build_hosts_candidate_datasets_table.py
\begin{landscape}
\begingroup
\footnotesize
\setlength{\tabcolsep}{4pt}
\renewcommand{\arraystretch}{1.15}
\begin{longtable}{@{}r@{\hspace{4pt}} L{78mm} L{72mm} L{76mm}@{}}
\caption[]{Sources for the recommended fifty and the ten extra of Table~\ref{tab:final}, in the same order.
\emph{Why it is ordered here} is the band reason for a host-bootstrap, regulatory-DNA or
cross-host-fitness row, and the tier and measurement rule for a scale row. \emph{Where
the per-record data lives} is the deposit, with its retrieval class: a plain entry is
fetched programmatically, \emph{browser} needs a human because the host returns 403 or
renders in JavaScript, and \emph{ask} means no public copy exists and the request is in
Table~\ref{tab:asks}. Every link is clickable.}
\label{tab:sources}\\
\toprule
\textbf{\#} & \textbf{Dataset, citation and link} & \textbf{Why it is ordered here} & \textbf{Where the per-record data lives} \\
\midrule
\endfirsthead
\multicolumn{4}{@{}l}{\footnotesize\emph{Table~\ref{tab:sources}, continued}}\\
\toprule
\textbf{\#} & \textbf{Dataset, citation and link} & \textbf{Why it is ordered here} & \textbf{Where the per-record data lives} \\
\midrule
\endhead
\bottomrule
\endfoot
1 & \textbf{Douglass 2018 (32 sequenced isolates, 4-antifungal MIC)} \org{P. kudriavzevii}\newline Douglass AP, Offei B, Braun-Galleani S, Coughlan AY, Martos AAR, Ortiz-Merino RA, Byrne KP, Wolfe KH. PLoS Pathog 2018;14:e1007138.\newline \href{https://www.ncbi.nlm.nih.gov/pmc/articles/PMC6053246/}{PMC60532\allowbreak 46} & Owner pin, 2026-11-24, to rank 1 from the host bootstrap band: it is the one P. kudriavzevii growth-type panel that can be fetched today with no ask, and fluconazole is the one compound measured on natural isolates in both yeasts, against Peter 2018 at 20 ug/ml on 971 S. cerevisiae isolates and Bloom 2019 at 100 uM on 13,950 segregants, so it is the first cross-host transfer test that has labels at both ends. The other three drugs are measured in S. cerevisiae only as deletion or CRISPRi chemogenomics, in Parsons 2006, Hillenmeyer 2008, and Piotrowski 2017. & SRA SRP1\allowbreak 39299 (3\allowbreak 2 strain\allowbreak s); GenB\allowbreak ank CP02\allowbreak 8773-\allowbreak CP02877\allowbreak 8 (CBS57\allowbreak 3 refere\allowbreak nce) and\allowbreak  CP02853\allowbreak 1-\allowbreak CP02853\allowbreak 5 (CBS51\allowbreak 47); MIC\allowbreak  table i\allowbreak n the op\allowbreak en-\allowbreak access \allowbreak article \allowbreak body \\
\addlinespace[5pt]
2 & \textbf{de Boer 2020 (GPRA sort-seq, 100M random promoters)} \org{S. cerevisiae}\newline de Boer CG, Vaishnav ED, Sadeh R, Abeyta EL, Friedman N, Regev A. Nat Biotechnol 2020;38:56-65.\newline \href{https://doi.org/10.1038/s41587-019-0315-8}{doi:\allowbreak 10.\allowbreak 1038/\allowbreak s41587-\allowbreak 019-\allowbreak 0315-\allowbreak 8} & Regulatory DNA band: one row carrying more measurements than every other row in this wave combined, on a schema that already serves its child, so the admission is a download decision rather than a design one. & GEO GSE1\allowbreak 04878 an\allowbreak d GSE104\allowbreak 903; Add\allowbreak gene 127\allowbreak 546; whi\allowbreak ch acces\allowbreak sion hol\allowbreak ds which\allowbreak  library\allowbreak  is unre\allowbreak solved a\allowbreak nd one p\allowbreak age chec\allowbreak k settle\allowbreak s it \\
\addlinespace[5pt]
3 & \textbf{Vaishnav 2022 deferred training libraries} \org{S. cerevisiae}\newline Vaishnav ED, de Boer CG, Molinet J, Yassour M, Fan L, Adiconis X, Thompson DA, Levin JZ, Cubillos FA, Regev A. Nature 2022;603:455-463.\newline \href{https://doi.org/10.1038/s41586-022-04506-6}{doi:\allowbreak 10.\allowbreak 1038/\allowbreak s41586-\allowbreak 022-\allowbreak 04506-\allowbreak 6} & Regulatory DNA band, and the single lowest-effort large win in the wave: every piece of machinery it needs is already written and serving its subset. & Zenodo 4\allowbreak 436477 (\allowbreak complex\_\allowbreak media\_\allowbreak trainin\allowbreak g\_\allowbreak data\_\allowbreak Glu.\allowbreak txt 3.\allowbreak 59 GB, \allowbreak defined\_\allowbreak media\_\allowbreak trainin\allowbreak g\_\allowbreak data\_\allowbreak SC\_\allowbreak Ura.\allowbreak txt 2.\allowbreak 47 GB);\allowbreak  GEO GSE\allowbreak 163045 a\allowbreak nd GSE16\allowbreak 3866 \\
\addlinespace[5pt]
4 & \textbf{Hollerer 2020 (ribosome binding site library, uASPIre)} \org{E. coli}\newline Hollerer S, Papaxanthos L, Gumpinger AC, Fischer K, Beisel C, Borgwardt K, Benenson Y, Jeschek M. Large-scale DNA-based phenotypic recording and deep learning enable highly accurate sequence-function mapping. Nat Commun 2020;11:3551.\newline \href{https://doi.org/10.1038/s41467-020-17222-4}{doi:\allowbreak 10.\allowbreak 1038/\allowbreak s41467-\allowbreak 020-\allowbreak 17222-\allowbreak 4} & Regulatory DNA: a 300,000-member binding-site library whose readout is a DNA-recorded time integral, a measurement principle no other row in the band carries. & the Jesc\allowbreak hekLab/\allowbreak uASPIre\allowbreak  reposit\allowbreak ory, rea\allowbreak chable \\
\addlinespace[5pt]
5 & \textbf{Hollerer 2023 (5' UTR and CDS library, uASPIre)} \org{E. coli}\newline Hollerer S, Jeschek M. Nucleic Acids Res 2023;51:2377-2396.\newline \href{https://doi.org/10.1093/nar/gkad040}{doi:\allowbreak 10.\allowbreak 1093/\allowbreak nar/\allowbreak gkad040} & Regulatory DNA: the band's whole purpose, an axis that adds measurements rather than rows, and this row adds more of them than any other non-yeast candidate. & Zenodo 1\allowbreak 0.\allowbreak 5281/\allowbreak zenodo.\allowbreak 7536586\allowbreak  and the\allowbreak  Jeschek\allowbreak Lab/\allowbreak uASPIre\_\allowbreak UTR\_\allowbreak CDS rep\allowbreak ository,\allowbreak  both re\allowbreak achable \\
\addlinespace[5pt]
6 & \textbf{DREAM 2022 random promoter challenge (Rafi 2024)} \org{S. cerevisiae}\newline Rafi AM, et al. A community effort to optimize sequence-based deep learning models of gene regulation. Nat Biotechnol 2024.\newline \href{https://www.ncbi.nlm.nih.gov/geo/query/acc.cgi?acc=GSE254493}{ncbi.\allowbreak nlm.\allowbreak nih.\allowbreak gov/\allowbreak geo/\allowbreak query/\allowbreak acc.\allowbreak cgi?\allowbreak acc=\allowbreak GSE2544\allowbreak 93} & Regulatory DNA band: the benchmark slice, which is what makes a model comparison possible at all, and it carries its own eight-category test design. & GEO GSE2\allowbreak 54493; Z\allowbreak enodo 10\allowbreak 633252 (\allowbreak v2, abou\allowbreak t 4.\allowbreak 0 GB), \allowbreak supersed\allowbreak ing Zeno\allowbreak do 73953\allowbreak 97 \\
\addlinespace[5pt]
7 & \textbf{Kuo 2020 (Shine-Dalgarno 9-mer library in four contexts)} \org{E. coli}\newline Kuo ST, Jahn RL, Cheng YJ, Chen YH, Lee YJ, Hollfelder F, Wen JD, Chou HHD. Global fitness landscapes of the Shine-Dalgarno sequence. Genome Res 2020;30:711-723.\newline \href{https://doi.org/10.1101/gr.260182.119}{doi:\allowbreak 10.\allowbreak 1101/\allowbreak gr.\allowbreak 260182.\allowbreak 119} & Regulatory DNA: a complete nine-mer space in four contexts, which is the band's axis measured exhaustively rather than sampled. & Europe P\allowbreak MC suppl\allowbreak ementary\allowbreak  files u\allowbreak nder PMC\allowbreak 7263185,\allowbreak  48,362,\allowbreak 217 byte\allowbreak s; reads\allowbreak  at PRJN\allowbreak A516114;\allowbreak  the pub\allowbreak lisher's\allowbreak  own sit\allowbreak e return\allowbreak s 403 \\
\addlinespace[5pt]
8 & \textbf{Kotopka and Smolke 2020 (artificial promoter libraries)} \org{S. cerevisiae}\newline Kotopka BJ, Smolke CD. Model-driven generation of artificial yeast promoters. Nat Commun 2020;11:2113.\newline \href{https://doi.org/10.1038/s41467-020-15977-4}{doi:\allowbreak 10.\allowbreak 1038/\allowbreak s41467-\allowbreak 020-\allowbreak 15977-\allowbreak 4} & Regulatory DNA band: the only million-scale promoter library outside the GPRA lineage, which is what makes it a check on it rather than more of it. & Zenodo 3\allowbreak 735426 (\allowbreak Promoter\allowbreak  Design \allowbreak Data\_\allowbreak 2020-\allowbreak 03.\allowbreak zip, ab\allowbreak out 1.\allowbreak 1 GB) \\
\addlinespace[5pt]
9 & \textbf{Zhu 2024 (Pol II MASTER, TSS initiation efficiency)} \org{S. cerevisiae}\newline Zhu C, Vvedenskaya IO, Sze SH, Nickels BE, Kaplan CD. Nat Struct Mol Biol 2024;31:190-202.\newline \href{https://doi.org/10.1038/s41594-023-01171-9}{doi:\allowbreak 10.\allowbreak 1038/\allowbreak s41594-\allowbreak 023-\allowbreak 01171-\allowbreak 9} & Regulatory DNA band: a different readout from reporter level, so it extends the axis rather than deepening it, and at a million records it is the third-largest row in the wave. & GEO GSE1\allowbreak 85290 (p\allowbreak rocessed\allowbreak ); BioPr\allowbreak oject PR\allowbreak JNA76662\allowbreak 4 (raw) \\
\addlinespace[5pt]
10 & \textbf{Savinov 2021 (random 50-mer 3' UTR library)} \org{S. cerevisiae}\newline Savinov A, Brandsen BM, Angell BE, Cuperus JT, Fields S. Genome Biol 2021;22:293.\newline \href{https://doi.org/10.1186/s13059-021-02509-6}{doi:\allowbreak 10.\allowbreak 1186/\allowbreak s13059-\allowbreak 021-\allowbreak 02509-\allowbreak 6} & Regulatory DNA: the 3' UTR end of the transcript, the one regulatory axis no other row in the band measures. & a figsha\allowbreak re recor\allowbreak d servin\allowbreak g a 36,9\allowbreak 20,254-\allowbreak byte ar\allowbreak chive th\allowbreak rough it\allowbreak s API; r\allowbreak eads at \allowbreak SRA PRJN\allowbreak A750726 \\
\addlinespace[5pt]
11 & \textbf{Urtecho 2020 (genome-wide promoter MPRA)} \org{E. coli}\newline Urtecho G, et al. Genome-wide functional characterization of Escherichia coli promoters and regulatory elements responsible for their function. eLife Reviewed Preprint 2023;12:92558.\newline \href{https://doi.org/10.7554/eLife.92558.1}{doi:\allowbreak 10.\allowbreak 7554/\allowbreak eLife.\allowbreak 92558.\allowbreak 1} & Regulatory DNA: the only row that tiles a whole bacterial genome for promoter activity, so its sequences are native rather than designed. & GEO GSE1\allowbreak 44621, r\allowbreak eachable\allowbreak  over FT\allowbreak P; the a\allowbreak ccession\allowbreak  display\allowbreak  page is\allowbreak  behind \allowbreak a challe\allowbreak nge and \allowbreak the auth\allowbreak ors' own\allowbreak  databas\allowbreak e fails \allowbreak TLS \\
\addlinespace[5pt]
12 & \textbf{Cuperus 2017 (random 5' UTR library)} \org{S. cerevisiae}\newline Cuperus JT, Groves B, Kuchina A, Rosenberg AB, Jojic N, Fields S, Seelig G. Genome Res 2017;27:2015-2028.\newline \href{https://genome.cshlp.org/content/27/12/2015}{genome.\allowbreak cshlp.\allowbreak org/\allowbreak content\allowbreak /\allowbreak 27/\allowbreak 12/\allowbreak 2015} & Regulatory DNA band: the only large 5' UTR library, and the only one in the band whose readout is growth rather than fluorescence. & GEO GSE1\allowbreak 04252 \\
\addlinespace[5pt]
13 & \textbf{Johns 2018 (29,249 expression constructs across conditions)} \org{E. coli}\newline Johns NI, et al. Metagenomic mining of regulatory elements enables programmable species-selective gene expression. Nat Methods 2018;15:323-329.\newline \href{https://doi.org/10.1038/nmeth.4633}{doi:\allowbreak 10.\allowbreak 1038/\allowbreak nmeth.\allowbreak 4633} & Regulatory DNA: the only library in the band whose sequences were mined from metagenomes rather than designed or tiled from one genome. & two publ\allowbreak isher su\allowbreak pplement\allowbreak ary file\allowbreak s, both \allowbreak reachabl\allowbreak e; no GE\allowbreak O or SRA\allowbreak  accessi\allowbreak on found\allowbreak , and th\allowbreak e full-\allowbreak text se\allowbreak rvice re\allowbreak turns an\allowbreak  error o\allowbreak n the av\allowbreak ailabili\allowbreak ty state\allowbreak ment \\
\addlinespace[5pt]
14 & \textbf{Yu 2021 (promoter library across induction states)} \org{E. coli}\newline Yu H, et al. Dynamic and quantitative characterization of promoters in Escherichia coli. Nat Commun 2021;12:325.\newline \href{https://doi.org/10.1038/s41467-020-20094-3}{doi:\allowbreak 10.\allowbreak 1038/\allowbreak s41467-\allowbreak 020-\allowbreak 20094-\allowbreak 3} & Regulatory DNA: a chromosomal promoter library carrying an induction response per sequence rather than a single level. & GEO GSE1\allowbreak 45630, r\allowbreak eachable\allowbreak  over FT\allowbreak P \\
\addlinespace[5pt]
15 & \textbf{Cambray 2018 (244,000 designed sequences, multi-phenotype)} \org{E. coli}\newline Cambray G, Guimaraes JC, Arkin AP. Evaluation of 244,000 synthetic sequences reveals design principles to optimize translation in Escherichia coli. Nat Biotechnol 2018;36:1005-1015.\newline \href{https://doi.org/10.1038/nbt.4238}{doi:\allowbreak 10.\allowbreak 1038/\allowbreak nbt.\allowbreak 4238} & Regulatory DNA: one designed library of 244,000 sequences read out five different ways, the widest phenotype fan-out in the wave. & an OSF d\allowbreak eposit a\allowbreak t 130,12\allowbreak 2,456 by\allowbreak tes that\allowbreak  answers\allowbreak  a scrip\allowbreak t, where\allowbreak  the pub\allowbreak lisher's\allowbreak  own lin\allowbreak k redire\allowbreak cts; rea\allowbreak ds at SR\allowbreak P086076 \\
\addlinespace[5pt]
16 & \textbf{Fuqua 2024 (promoter-island mutational landscape)} \org{E. coli}\newline Fuqua T, et al. Dense and pleiotropic regulatory information in a developmental enhancer-like bacterial promoter island. eLife 2024;13:98654.\newline \href{https://doi.org/10.7554/eLife.98654}{doi:\allowbreak 10.\allowbreak 7554/\allowbreak eLife.\allowbreak 98654} & Regulatory DNA: a quarter-million mutated native promoter sequences released as per-bin distributions rather than means. & the tfuq\allowbreak ua95/\allowbreak promote\allowbreak r\_\allowbreak islands\allowbreak  reposit\allowbreak ory, rea\allowbreak chable; \allowbreak reads at\allowbreak  SRA PRJ\allowbreak NA107157\allowbreak 2 \\
\addlinespace[5pt]
17 & \textbf{Gui 2021 (position effects with transcriptomes)} \org{S. cerevisiae}\newline Gui, et al. Transcriptome analysis in yeast reveals the externality of position effects. Mol Biol Evol 2021.\newline \href{https://www.biorxiv.org/content/10.1101/2020.04.02.021162}{biorxiv.\allowbreak org/\allowbreak content\allowbreak /\allowbreak 10.\allowbreak 1101/\allowbreak 2020.\allowbreak 04.\allowbreak 02.\allowbreak 021162} & Regulatory DNA band: the sequence variable is purely the integration coordinate, and this is the one row in that tier whose phenotype is a vector. & no acces\allowbreak sion fou\allowbreak nd; need\allowbreak s a brow\allowbreak ser chec\allowbreak k of the\allowbreak  publish\allowbreak ed suppl\allowbreak ement\newline {\scriptsize (browser)} \\
\addlinespace[5pt]
18 & \textbf{Zhou 2025 (massively parallel bacterial promoter library)} \org{E. coli}\newline Zhou Y, et al. Nucleic Acids Res 2025;53:gkaf863.\newline \href{https://doi.org/10.1093/nar/gkaf863}{doi:\allowbreak 10.\allowbreak 1093/\allowbreak nar/\allowbreak gkaf863} & Regulatory DNA: 700,407 promoter measurements, and the clearest case in the wave of a release exceeding its own paper's claim. & Europe P\allowbreak MC suppl\allowbreak ementary\allowbreak  files u\allowbreak nder PMC\allowbreak 12407095\allowbreak ; reads \allowbreak at PRJNA\allowbreak 1189185 \\
\addlinespace[5pt]
19 & \textbf{Keren 2016 (promoter panel x named gene, pooled fitness)} \org{S. cerevisiae}\newline Keren L, et al. Massively parallel interrogation of the effects of gene expression levels on fitness. Cell 2016;166:1282-1294.\newline \href{https://www.ncbi.nlm.nih.gov/geo/query/acc.cgi?acc=GSE83936}{ncbi.\allowbreak nlm.\allowbreak nih.\allowbreak gov/\allowbreak geo/\allowbreak query/\allowbreak acc.\allowbreak cgi?\allowbreak acc=\allowbreak GSE8393\allowbreak 6} & Regulatory DNA band: the only large tabulated panel that links expression level to fitness gene by gene, and the only one whose perturbation is a swap at a native locus. & GEO GSE8\allowbreak 3936; Ar\allowbreak rayExpre\allowbreak ss E-\allowbreak GEOD-\allowbreak 83936; \allowbreak BioProje\allowbreak ct PRJNA\allowbreak 329129 \\
\addlinespace[5pt]
20 & \textbf{Shalem 2015 (3' end FACS-seq library)} \org{S. cerevisiae}\newline Shalem O, et al. 3' end determinants. PLoS Genet 2015;11:e1005147.\newline \href{https://doi.org/10.1371/journal.pgen.1005147}{doi:\allowbreak 10.\allowbreak 1371/\allowbreak journal\allowbreak .\allowbreak pgen.\allowbreak 1005147} & Regulatory DNA band: the terminator end of the same axis, with the sequence and the value in one released file. & PLoS Gen\allowbreak et S3 Ta\allowbreak ble (XLS\allowbreak X), no a\allowbreak ccession\allowbreak ; needs \allowbreak a supple\allowbreak ment fet\allowbreak ch and a\allowbreak  sha256 \allowbreak pin\newline {\scriptsize (browser)} \\
\addlinespace[5pt]
21 & \textbf{Renganaath 2020 (native promoter variant MPRA)} \org{S. cerevisiae}\newline Renganaath K, Cheung R, Day L, Kosuri S, Kruglyak L, Albert FW. eLife 2020;9:e62669.\newline \href{https://doi.org/10.7554/eLife.62669}{doi:\allowbreak 10.\allowbreak 7554/\allowbreak eLife.\allowbreak 62669} & Regulatory DNA band: real alleles rather than sampled sequence, measured on the same readout, which is the bridge between the synthetic and the natural axes. & GEO GSE1\allowbreak 55944 (r\allowbreak aw plus \allowbreak barcode \allowbreak assignme\allowbreak nts); th\allowbreak e MPRA d\allowbreak esign an\allowbreak d proces\allowbreak sed tabl\allowbreak es are e\allowbreak Life Sup\allowbreak plementa\allowbreak ry Files \\
\addlinespace[5pt]
22 & \textbf{Arg-CCU tRNA deep mutational scan} \org{S. cerevisiae}\newline Deep mutational scan of a single-copy Arg-CCU tRNA over 14 substitutions seen across post-whole-genome-duplication Saccharomyces. PMC6193533.\newline \href{https://pmc.ncbi.nlm.nih.gov/articles/PMC6193533}{PMC61935\allowbreak 33} & Regulatory DNA band: point substitutions in a non-coding gene, which is the one DMS shape the protein aggregates cannot hold. & journal \allowbreak suppleme\allowbreak nt, no a\allowbreak ccession\allowbreak  located\allowbreak ; needs \allowbreak a supple\allowbreak ment fet\allowbreak ch\newline {\scriptsize (browser)} \\
\addlinespace[5pt]
23 & \textbf{Keren 2013 (native promoter panel across conditions)} \org{S. cerevisiae}\newline Keren L, et al. Mol Syst Biol 2013.\newline \href{https://spell.yeastgenome.org/}{spell.\allowbreak yeastge\allowbreak nome.\allowbreak org} & Regulatory DNA band: the only row in it with a real condition axis on native sequence. & not loca\allowbreak ted; rea\allowbreak ched onl\allowbreak y throug\allowbreak h the Zr\allowbreak imec and\allowbreak  Karollu\allowbreak s valida\allowbreak tion set\allowbreak s\newline {\scriptsize (browser)} \\
\addlinespace[5pt]
24 & \textbf{Lubliner 2015 (core promoter variant panel)} \org{S. cerevisiae}\newline Lubliner S, Keren L, Segal E, et al. Core promoter sequence in yeast is a major determinant of expression level. Genome Res 2015;25:1008-1017.\newline \href{https://genome.cshlp.org/content/25/7/1008}{genome.\allowbreak cshlp.\allowbreak org/\allowbreak content\allowbreak /\allowbreak 25/\allowbreak 7/\allowbreak 1008} & Regulatory DNA band: the core promoter is a distinct regulatory element from the upstream activating region every other row varies. & Genome R\allowbreak es Suppl\allowbreak emental \allowbreak Table 1;\allowbreak  needs a\allowbreak  supplem\allowbreak ent fetc\allowbreak h\newline {\scriptsize (browser)} \\
\addlinespace[5pt]
25 & \textbf{Sharon 2012 (designed promoter library)} \org{S. cerevisiae}\newline Sharon E, Kalma Y, Sharp A, et al., Segal E. Nat Biotechnol 2012;30:521-530.\newline \href{https://doi.org/10.1038/nbt.2205}{doi:\allowbreak 10.\allowbreak 1038/\allowbreak nbt.\allowbreak 2205} & Regulatory DNA band: the controlled counterpart the random libraries lack, and the platform's own origin. & Nat Biot\allowbreak echnol j\allowbreak ournal s\allowbreak upplemen\allowbreak t, no ac\allowbreak cession;\allowbreak  needs a\allowbreak  supplem\allowbreak ent fetc\allowbreak h\newline {\scriptsize (browser)} \\
\addlinespace[5pt]
26 & \textbf{Yamanishi 2013 (genomic terminatome)} \org{S. cerevisiae}\newline Yamanishi, et al. Characterization of heterologous protein expression using over 3,000 yeast genomic terminators. SGD reference S000206762.\newline \href{https://www.yeastgenome.org/reference/S000206762}{yeastgen\allowbreak ome.\allowbreak org/\allowbreak referen\allowbreak ce/\allowbreak S000206\allowbreak 762} & Regulatory DNA band: the largest native-terminator panel, which is the one element of the axis no library in the band randomizes. & primary \allowbreak table no\allowbreak t locate\allowbreak d; needs\allowbreak  the pap\allowbreak er ident\allowbreak ified, t\allowbreak hen a su\allowbreak pplement\allowbreak  fetch\newline {\scriptsize (browser)} \\
\addlinespace[5pt]
27 & \textbf{Dvir 2013 (5' UTR positions -10 to -1)} \org{S. cerevisiae}\newline Dvir S, Velten L, Sharon E, Zeevi D, Carey LB, Weinberger A, Segal E. Proc Natl Acad Sci USA 2013;110:E2792-E2801.\newline \href{https://doi.org/10.1073/pnas.1222534110}{doi:\allowbreak 10.\allowbreak 1073/\allowbreak pnas.\allowbreak 1222534\allowbreak 110} & Regulatory DNA band: a controlled ten-base window, which is the designed counterpart to the half-million random 5' UTRs above it. & PNAS jou\allowbreak rnal sup\allowbreak plement,\allowbreak  no acce\allowbreak ssion; n\allowbreak eeds a s\allowbreak upplemen\allowbreak t fetch\newline {\scriptsize (browser)} \\
\addlinespace[5pt]
28 & \textbf{Wu 2017 (genome-wide position effects, 1,044 loci)} \org{S. cerevisiae}\newline Wu XL, et al. Genome-wide landscape of position effects on heterogeneous gene expression in Saccharomyces cerevisiae. Biotechnol Biofuels 2017.\newline \href{https://www.ncbi.nlm.nih.gov/pmc/articles/PMC5516366/}{PMC55163\allowbreak 66} & Regulatory DNA band by its sequence variable, the integration coordinate, even though the readout is one reporter value. & Table S2\allowbreak , stated\allowbreak  availab\allowbreak le from \allowbreak the corr\allowbreak espondin\allowbreak g author\allowbreak  on reas\allowbreak onable r\allowbreak equest\newline {\scriptsize (browser)} \\
\addlinespace[5pt]
29 & \textbf{ProteinGym yeast subset} \org{S. cerevisiae}\newline ProteinGym: deep mutational scanning scores from 217 assays over 186 proteins, distributed through ProteinGymR and ProteinGymPy.\newline \href{https://proteingym.org}{proteing\allowbreak ym.\allowbreak org} & Regulatory DNA band as the aggregate route into the mutational-scan tier, which is the sequence-to-function axis in coding space rather than regulatory space. & ProteinG\allowbreak ymR and \allowbreak ProteinG\allowbreak ymPy pac\allowbreak kage loa\allowbreak ders; ar\allowbreak chived o\allowbreak n Zenodo\allowbreak , record\allowbreak  number \allowbreak not capt\allowbreak ured \\
\addlinespace[5pt]
30 & \textbf{Zhang 2016 (genome-wide GFP position scan)} \org{S. cerevisiae}\newline Zhang, et al. Genome-wide GFP position-effect scan using the yeast knockout collection. Chin J Biotech 2016.\newline \href{https://wsws.ijournals.cn/cjbcn/article/pdf/gc16070901}{wsws.\allowbreak ijourna\allowbreak ls.\allowbreak cn/\allowbreak cjbcn/\allowbreak article\allowbreak /\allowbreak pdf/\allowbreak gc16070\allowbreak 901} & Regulatory DNA band by its sequence variable, and the only FACS-read position-effect scan. & no acces\allowbreak sion loc\allowbreak ated; Ch\allowbreak inese-\allowbreak languag\allowbreak e journa\allowbreak l PDF, n\allowbreak eeds a b\allowbreak rowser f\allowbreak etch\newline {\scriptsize (browser)} \\
\addlinespace[5pt]
31 & \textbf{Kita 2017 (context-stratified yeast eQTLs)} \org{S. cerevisiae}\newline Kita R, et al. Proc Natl Acad Sci USA 2017;114.\newline \href{https://www.pnas.org/doi/10.1073/pnas.1717421114}{pnas.\allowbreak org/\allowbreak doi/\allowbreak 10.\allowbreak 1073/\allowbreak pnas.\allowbreak 1717421\allowbreak 114} & Regulatory DNA band: variants stratified by the regulatory element they sit in, which is the only released set with that stratification. & PNAS sup\allowbreak plementa\allowbreak ry data \allowbreak file pna\allowbreak s.\allowbreak 1717421\allowbreak 114.\allowbreak sd01; n\allowbreak eeds a s\allowbreak upplemen\allowbreak t fetch\newline {\scriptsize (browser)} \\
\addlinespace[5pt]
32 & \textbf{Chen (earlier position-effect study, 482 loci)} \org{S. cerevisiae}\newline An earlier Chen study characterizing position effects at 482 loci, reported inside Wu 2017; primary citation not located.\newline \href{https://www.ncbi.nlm.nih.gov/pmc/articles/PMC5516366/}{PMC55163\allowbreak 66} & Regulatory DNA band by its sequence variable, the integration coordinate. & not loca\allowbreak ted; the\allowbreak  primary\allowbreak  paper m\allowbreak ust be i\allowbreak dentifie\allowbreak d first\newline {\scriptsize (browser)} \\
\addlinespace[5pt]
33 & \textbf{Wei 2017 (synthetic terminator panel)} \org{S. cerevisiae}\newline Wei, et al. Yeast synthetic terminators: fine regulation of strength through linker sequences. ChemBioChem 2017.\newline \href{https://doi.org/10.1002/cbic.201700516}{doi:\allowbreak 10.\allowbreak 1002/\allowbreak cbic.\allowbreak 2017005\allowbreak 16} & Regulatory DNA band: the terminator parts tier, whose values are relative to a standard part rather than absolute. & ChemBioC\allowbreak hem jour\allowbreak nal supp\allowbreak lement, \allowbreak no repos\allowbreak itory de\allowbreak posit fo\allowbreak und\newline {\scriptsize (browser)} \\
\addlinespace[5pt]
34 & \textbf{Thompson 2019 (KT2440 lysine metabolism)} \org{P. putida}\newline Thompson MG, et al. Massively parallel fitness profiling reveals multiple novel enzymes in Pseudomonas putida lysine metabolism. mBio 2019.\newline \href{https://pmc.ncbi.nlm.nih.gov/articles/PMC6509195}{PMC65091\allowbreak 95} & Scale band: tier 3, ordered by measurements (22,256). & Table S1\allowbreak , a CSV \allowbreak suppleme\allowbreak ntary fi\allowbreak le; no r\allowbreak epositor\allowbreak y\newline {\scriptsize (browser)} \\
\addlinespace[5pt]
35 & \textbf{Morse (short synthetic terminators)} \org{S. cerevisiae}\newline Morse NJ, et al. Short synthetic terminators, 35 to 70 bp; primary citation not located.\newline \href{https://scite.ai/authors/nicholas-j-morse-1g5Qrn}{scite.\allowbreak ai/\allowbreak authors\allowbreak /\allowbreak nichola\allowbreak s-\allowbreak j-\allowbreak morse-\allowbreak 1g5Qrn} & Regulatory DNA band: the short-terminator corner of the parts tier, held so the lead is not lost. & no repos\allowbreak itory de\allowbreak posit fo\allowbreak und; pri\allowbreak mary pap\allowbreak er not i\allowbreak dentifie\allowbreak d\newline {\scriptsize (browser)} \\
\addlinespace[5pt]
36 & \textbf{MaveDB yeast score sets} \org{S. cerevisiae}\newline Esposito D, et al. MaveDB: an open-source platform to distribute and interpret data from multiplexed assays of variant effect. Genome Biol 2019;20:223.\newline \href{https://doi.org/10.1186/s13059-019-1845-6}{doi:\allowbreak 10.\allowbreak 1186/\allowbreak s13059-\allowbreak 019-\allowbreak 1845-\allowbreak 6} & Regulatory DNA band as the second aggregate route into the mutational-scan tier, held because the two aggregates do not hold the same assays. & mavedb.\allowbreak org; th\allowbreak e docume\allowbreak nted sea\allowbreak rch API \allowbreak endpoint\allowbreak  returns\allowbreak  HTTP 40\allowbreak 4, so th\allowbreak e organi\allowbreak sm filte\allowbreak r needs \allowbreak a live q\allowbreak uery\newline {\scriptsize (browser)} \\
\addlinespace[5pt]
37 & \textbf{Isoprenol CRISPRi DBTL proteome ladder (7 PRIDE deposits)} \org{P. putida}\newline Automation of a CRISPRi platform for enhanced isoprenol production in Pseudomonas putida; PRIDE projects PXD063733, PXD063737, PXD063738, PXD063740, PXD063743, PXD063744, PXD063746.\newline \href{https://www.ebi.ac.uk/pride/archive/projects/PXD063733}{ebi.\allowbreak ac.\allowbreak uk/\allowbreak pride/\allowbreak archive\allowbreak /\allowbreak project\allowbreak s/\allowbreak PXD0637\allowbreak 33} & Scale band: tier 3, ordered by measurements (38,948). & PRIDE PX\allowbreak D063733,\allowbreak  PXD0637\allowbreak 37, PXD0\allowbreak 63738, P\allowbreak XD063740\allowbreak , PXD063\allowbreak 743, PXD\allowbreak 063744 a\allowbreak nd PXD06\allowbreak 3746, al\allowbreak l submit\allowbreak ted 2025\allowbreak -\allowbreak 05-\allowbreak 07 and \allowbreak 08 on an\allowbreak  Orbitra\allowbreak p Explor\allowbreak is 480 \\
\addlinespace[5pt]
38 & \textbf{Thompson 2020 (KT2440 fatty acids and alcohols)} \org{P. putida}\newline Thompson MG, Incha MR, Pearson AN, et al. Appl Environ Microbiol 2020;86:e01665-20; correction at 10.1128/AEM.00177-21.\newline \href{https://pubmed.ncbi.nlm.nih.gov/32826213}{PMID 328\allowbreak 26213} & Scale band: tier 3, ordered by measurements (127,972). & not stat\allowbreak ed; the \allowbreak values m\allowbreak ost like\allowbreak ly reach\allowbreak  feba.\allowbreak db \\
\addlinespace[5pt]
39 & \textbf{GSE198395 (putidaPRECISE321 carbon sources)} \org{P. putida}\newline GEO series GSE198395, KT2440 and engineered strains on different carbon sources, part of putidaPRECISE321, released 2022-04-28.\newline \href{https://www.ncbi.nlm.nih.gov/geo/query/acc.cgi?acc=GSE198395}{ncbi.\allowbreak nlm.\allowbreak nih.\allowbreak gov/\allowbreak geo/\allowbreak query/\allowbreak acc.\allowbreak cgi?\allowbreak acc=\allowbreak GSE1983\allowbreak 95} & Scale band: tier 3, ordered by measurements (450,684). & GEO GSE1\allowbreak 98395; p\allowbreak rogramma\allowbreak tically \allowbreak retrieva\allowbreak ble \\
\addlinespace[5pt]
40 & \textbf{Mueller 2025 (24-strain P. putida Biolog panel)} \org{P. putida}\newline Mueller J, Krishnan J, Wei Q, Hefner Y, Monk JM, Verkler H, Tibocha-Bonilla JD, et al. bioRxiv 2025;2025.11.17.688870.\newline \href{https://doi.org/10.1101/2025.11.17.688870}{doi:\allowbreak 10.\allowbreak 1101/\allowbreak 2025.\allowbreak 11.\allowbreak 17.\allowbreak 688870} & Scale band: tier 2, ordered by measurements (4,920). & unconfir\allowbreak med:\allowbreak  bioRxi\allowbreak v return\allowbreak ed HTTP \allowbreak 429 on f\allowbreak our atte\allowbreak mpts, so\allowbreak  no BioP\allowbreak roject, \allowbreak Zenodo r\allowbreak ecord or\allowbreak  machine\allowbreak -\allowbreak readabl\allowbreak e Table \allowbreak S2 is co\allowbreak nfirmed\newline {\scriptsize (ask)} \\
\addlinespace[5pt]
41 & \textbf{Schmidt 2022 (KT2440 nitrogen BarSeq)} \org{P. putida}\newline Schmidt M, Pearson AN, Incha MR, et al. Appl Environ Microbiol 2022;88:e02430-21.\newline \href{https://pmc.ncbi.nlm.nih.gov/articles/PMC9004399}{PMC90043\allowbreak 99} & Scale band: tier 2, ordered by measurements (395,044). & no acces\allowbreak sion sta\allowbreak ted; the\allowbreak  study l\allowbreak inks int\allowbreak o the Fi\allowbreak tness Br\allowbreak owser, s\allowbreak o its ex\allowbreak periment\allowbreak s are mo\allowbreak st likel\allowbreak y inside\allowbreak  feba.\allowbreak db \\
\addlinespace[5pt]
42 & \textbf{Borchert 2024 (KT2440 fitness compendium)} \org{P. putida}\newline Borchert AJ, Bleem AC, Lim HG, Rychel K, Dooley KD, Kellermyer ZA, Hodges TL, Palsson BO, Beckham GT. mSystems 2024;9:e00942-23.\newline \href{https://doi.org/10.1128/msystems.00942-23}{doi:\allowbreak 10.\allowbreak 1128/\allowbreak msystem\allowbreak s.\allowbreak 00942-\allowbreak 23} & Scale band: tier 1, ordered by measurements (1,571,024). & github.\allowbreak com/\allowbreak beckham\allowbreak -\allowbreak lab/\allowbreak fModule\allowbreak  (fitnes\allowbreak s values\allowbreak , statis\allowbreak tics, me\allowbreak tadata);\allowbreak  the non\allowbreak -\allowbreak browser\allowbreak  samples\allowbreak  are at \allowbreak SRA PRJN\allowbreak A809672,\allowbreak  PRJNA85\allowbreak 6070, PR\allowbreak JNA10112\allowbreak 87 \\
\addlinespace[5pt]
43 & \textbf{putidaPRECISE321 (Lim 2022)} \org{P. putida}\newline Lim HG, Rychel K, Sastry AV, Mueller J, Niu W, Feist AM, Palsson BO. Metab Eng 2022;72.\newline \href{https://doi.org/10.1016/j.ymben.2022.04.004}{doi:\allowbreak 10.\allowbreak 1016/\allowbreak j.\allowbreak ymben.\allowbreak 2022.\allowbreak 04.\allowbreak 004} & Scale band: tier 1, ordered by measurements (1,786,044). & github.\allowbreak com/\allowbreak SBRG/\allowbreak modulom\allowbreak e\_\allowbreak ppu (MI\allowbreak T); name\allowbreak d GEO de\allowbreak posits G\allowbreak SE198395\allowbreak  (81 sam\allowbreak ples) an\allowbreak d GSE193\allowbreak 493 (22)\allowbreak ; the fu\allowbreak ll per-\allowbreak sample \allowbreak GEO list\allowbreak  is in t\allowbreak he SI sp\allowbreak readshee\allowbreak t \\
\addlinespace[5pt]
44 & \textbf{GSE262836 (Delhi NICU, fluconazole-resistance expression)} \org{P. kudriavzevii}\newline Genomic characterization and transcriptomic profile of Candida krusei in a neonatal intensive care unit in Delhi, India. GEO series GSE262836, 2025.\newline \href{https://www.ncbi.nlm.nih.gov/geo/query/acc.cgi?acc=GSE262836}{ncbi.\allowbreak nlm.\allowbreak nih.\allowbreak gov/\allowbreak geo/\allowbreak query/\allowbreak acc.\allowbreak cgi?\allowbreak acc=\allowbreak GSE2628\allowbreak 36} & Host bootstrap: one of two rows in this host pairing a drug phenotype with expression, and openly retrievable. & GEO GSE2\allowbreak 62836, p\allowbreak latform \allowbreak GPL34351\allowbreak ; progra\allowbreak mmatical\allowbreak ly retri\allowbreak evable \\
\addlinespace[5pt]
45 & \textbf{GSE159457 (CY902 C4 dicarboxylic acid transporters)} \org{P. kudriavzevii}\newline Identification and characterization of C4 dicarboxylic acid transporters in Pichia kudriavzevii CY902. GEO series GSE159457, 2021; PMID 33629807.\newline \href{https://www.ncbi.nlm.nih.gov/geo/query/acc.cgi?acc=GSE159457}{ncbi.\allowbreak nlm.\allowbreak nih.\allowbreak gov/\allowbreak geo/\allowbreak query/\allowbreak acc.\allowbreak cgi?\allowbreak acc=\allowbreak GSE1594\allowbreak 57} & Host bootstrap: a carbon-source axis on expression in a third background, fully deposited. & GEO GSE1\allowbreak 59457, p\allowbreak latform \allowbreak GPL29244\allowbreak ; raw re\allowbreak ads at S\allowbreak RA SRP28\allowbreak 7255; pr\allowbreak ogrammat\allowbreak ically r\allowbreak etrievab\allowbreak le \\
\addlinespace[5pt]
46 & \textbf{Dubinkina 2024 (12-strain low-pH transcriptomic atlas)} \org{P. kudriavzevii}\newline Dubinkina V, Bhogale S, Hsieh PH, Dibaeinia P, Nambiar S, Maslov S, Yoshikuni Y, Sinha S. Microbiol Spectr 2024.\newline \href{https://doi.org/10.1128/spectrum.02536-23}{doi:\allowbreak 10.\allowbreak 1128/\allowbreak spectru\allowbreak m.\allowbreak 02536-\allowbreak 23} & Host bootstrap: the only multi-strain expression matrix this host has, and the only deposit that already carries processed values rather than reads. & Illinois\allowbreak  Data Ba\allowbreak nk DOI 1\allowbreak 0.\allowbreak 13012/\allowbreak B2IDB-\allowbreak 1336177\_\allowbreak V1, hol\allowbreak ding str\allowbreak ain regu\allowbreak lators, \allowbreak strains \allowbreak and gene\allowbreak  sets, a\allowbreak nd gene \allowbreak expressi\allowbreak on count\allowbreak s; the p\allowbreak age retu\allowbreak rns 403 \allowbreak to every\allowbreak  automat\allowbreak ed clien\allowbreak t\newline {\scriptsize (browser)} \\
\addlinespace[5pt]
47 & \textbf{MAESTRO 2025 (170-isolate WGS, 57-stressor growth panel)} \org{P. kudriavzevii}\newline Hsieh PH, Sasaki Y, Yang CI, Wu ZY, Steindorff AS, Haridas S, Ke J, Fatma Z, Zhao Z, Opulente DA, Deng S, Hittinger CT, Grigoriev IV, Dien B, Zhao H, Li YP, Yoshikuni Y. bioRxiv 2025;2025.08.29.673034.\newline \href{https://doi.org/10.1101/2025.08.29.673034}{doi:\allowbreak 10.\allowbreak 1101/\allowbreak 2025.\allowbreak 08.\allowbreak 29.\allowbreak 673034} & Host bootstrap: this host has no genome-wide screen of any kind, so its largest resource is a 161-isolate panel that the measurement rule ranks below every pseudomonad screen. The band is what keeps it in the first tranche. & no acces\allowbreak sion ret\allowbreak rieved:\allowbreak  bioRxi\allowbreak v return\allowbreak s 403 an\allowbreak d Europe\allowbreak  PMC ful\allowbreak lTextXML\allowbreak  for PPR\allowbreak 1078985 \allowbreak returns \allowbreak 500; per\allowbreak -\allowbreak isolate\allowbreak  IO\_\allowbreak NNN Bio\allowbreak Projects\allowbreak  are inf\allowbreak erred, n\allowbreak ot confi\allowbreak rmed\newline {\scriptsize (ask)} \\
\addlinespace[5pt]
48 & \textbf{Fatma 2023 (landing-pad locus screen)} \org{P. kudriavzevii}\newline Fatma Z, Tan SI, Boob A, Zhao H. A landing pad system for multicopy gene integration in Issatchenkia orientalis. bioRxiv 2023;2023.05.21.541627.\newline \href{https://doi.org/10.1101/2023.05.21.541627}{doi:\allowbreak 10.\allowbreak 1101/\allowbreak 2023.\allowbreak 05.\allowbreak 21.\allowbreak 541627} & Host bootstrap: the host's only measured per-locus fitness, which no scale ranking would ever place in a first tranche at 188 records. & a 23 KB \allowbreak package \allowbreak at uofi.\allowbreak box.\allowbreak com (st\allowbreak rains, p\allowbreak lasmids,\allowbreak  guide R\allowbreak NAs, pri\allowbreak mers) pl\allowbreak us githu\allowbreak b.\allowbreak com/\allowbreak HuiminZ\allowbreak hao/\allowbreak Landing\allowbreak -\allowbreak pad-\allowbreak model; \allowbreak no DOI a\allowbreak nd no ar\allowbreak chival d\allowbreak eposit\newline {\scriptsize (ask)} \\
\addlinespace[5pt]
49 & \textbf{GSE282019 (Delhi NICU five-year outbreak)} \org{P. kudriavzevii}\newline Genomic analyses reveal high diversity and rapid evolution of Pichia kudriavzevii within a neonatal intensive care unit in Delhi, India. GEO series GSE282019, 2024; PMID 39853119.\newline \href{https://www.ncbi.nlm.nih.gov/geo/query/acc.cgi?acc=GSE282019}{ncbi.\allowbreak nlm.\allowbreak nih.\allowbreak gov/\allowbreak geo/\allowbreak query/\allowbreak acc.\allowbreak cgi?\allowbreak acc=\allowbreak GSE2820\allowbreak 19} & Host bootstrap: the host's only published drug-stratified expression series, and openly retrievable. & GEO GSE2\allowbreak 82019, p\allowbreak latform \allowbreak GPL34351\allowbreak ; progra\allowbreak mmatical\allowbreak ly retri\allowbreak evable \\
\addlinespace[5pt]
50 & \textbf{GSE202655 (L-malic acid, EcSthA transhydrogenase)} \org{P. kudriavzevii}\newline Metabolic engineering of the acid-tolerant yeast Pichia kudriavzevii for efficient L-malic acid production at low pH. GEO series GSE202655, 2023; PMID 36566973.\newline \href{https://www.ncbi.nlm.nih.gov/geo/query/acc.cgi?acc=GSE202655}{ncbi.\allowbreak nlm.\allowbreak nih.\allowbreak gov/\allowbreak geo/\allowbreak query/\allowbreak acc.\allowbreak cgi?\allowbreak acc=\allowbreak GSE2026\allowbreak 55} & Host bootstrap: an engineered-strain expression series in a non-SD108 background, which the host otherwise lacks. & GEO GSE2\allowbreak 02655, p\allowbreak latform \allowbreak GPL32243\allowbreak ; progra\allowbreak mmatical\allowbreak ly retri\allowbreak evable \\
\addlinespace[5pt]
51 & \textbf{Lubliner 2013 (native core promoter maximal strengths)} \org{S. cerevisiae}\newline Lubliner S, Keren L, Segal E. Nucleic Acids Res 2013.\newline \href{https://academic.oup.com/nar}{academic\allowbreak .\allowbreak oup.\allowbreak com/\allowbreak nar} & Regulatory DNA band: the native counterpart of the core-promoter tier, and the only source for its reference values. & not loca\allowbreak ted\newline {\scriptsize (browser)} \\
\addlinespace[5pt]
52 & \textbf{Fitness Browser RCH2 (Stutzerimonas stutzeri\_AE)} \org{S. stutzeri}\newline Price MN, Wetmore KM, Waters RJ, Callaghan M, Ray J, Liu H, Kuehl JV, Melnyk RA, Lamson JS, Suh Y, et al., Deutschbauer AM, Arkin AP. Nature 2018;557:503-509.\newline \href{https://doi.org/10.1038/s41586-018-0124-0}{doi:\allowbreak 10.\allowbreak 1038/\allowbreak s41586-\allowbreak 018-\allowbreak 0124-\allowbreak 0} & Cross-host fitness: nine new hosts arrive inside one artifact under one schema, one condition vocabulary and an identifier bridge the source supplies, so the band exists to keep them together rather than scattered down a measurement ranking. & feba.\allowbreak db at 1\allowbreak 0,433,70\allowbreak 8,032 by\allowbreak tes, Las\allowbreak t-\allowbreak Modifie\allowbreak d 2026-\allowbreak 10-\allowbreak 27, HTT\allowbreak P 200 to\allowbreak  a plain\allowbreak  curl; o\allowbreak rgId psR\allowbreak CH2; ref\allowbreak seq\_\allowbreak mapping\allowbreak .\allowbreak tsv bri\allowbreak dges loc\allowbreak usId to \allowbreak RefSeq \\
\addlinespace[5pt]
53 & \textbf{Fitness Browser FW300-N2E3 (Aquipseudomonas helvetica)} \org{P. fluorescens}\newline Price MN, Wetmore KM, Waters RJ, et al., Deutschbauer AM, Arkin AP. Nature 2018;557:503-509.\newline \href{https://doi.org/10.1038/s41586-018-0124-0}{doi:\allowbreak 10.\allowbreak 1038/\allowbreak s41586-\allowbreak 018-\allowbreak 0124-\allowbreak 0} & Cross-host fitness: part of the single-artifact block of nine hosts, whose shared schema is the reason to admit them together. & feba.\allowbreak db, org\allowbreak Id pseud\allowbreak o3\_\allowbreak N2E3; t\allowbreak he histo\allowbreak rical bi\allowbreak nomial i\allowbreak s P.\allowbreak  fluore\allowbreak scens an\allowbreak d the cu\allowbreak rrent on\allowbreak e is A.\allowbreak  helvet\allowbreak ica \\
\addlinespace[5pt]
54 & \textbf{Fitness Browser FW300-N2E2 (Aquipseudomonas fluorescens\_Q)} \org{P. fluorescens}\newline Price MN, Wetmore KM, Waters RJ, et al., Deutschbauer AM, Arkin AP. Nature 2018;557:503-509.\newline \href{https://doi.org/10.1038/s41586-018-0124-0}{doi:\allowbreak 10.\allowbreak 1038/\allowbreak s41586-\allowbreak 018-\allowbreak 0124-\allowbreak 0} & Cross-host fitness: one of four isolates of one species in the block, which is the within-species comparison the band makes possible. & feba.\allowbreak db, org\allowbreak Id pseud\allowbreak o6\_\allowbreak N2E2; 1\allowbreak 16 updat\allowbreak ed annot\allowbreak ations \\
\addlinespace[5pt]
55 & \textbf{Fitness Browser WCS417 (Aquipseudomonas simiae)} \org{P. simiae}\newline Price MN, Wetmore KM, Waters RJ, et al., Deutschbauer AM, Arkin AP. Nature 2018;557:503-509.\newline \href{https://doi.org/10.1038/s41586-018-0124-0}{doi:\allowbreak 10.\allowbreak 1038/\allowbreak s41586-\allowbreak 018-\allowbreak 0124-\allowbreak 0} & Cross-host fitness: the plant-associated member of the single-artifact block, with its widest carbon panel. & feba.\allowbreak db, org\allowbreak Id WCS41\allowbreak 7; histo\allowbreak rical bi\allowbreak nomial P\allowbreak .\allowbreak  simiae\allowbreak , curren\allowbreak t A.\allowbreak  simiae \\
\addlinespace[5pt]
56 & \textbf{Fitness Browser FW300-N2C3 (Aquipseudomonas fluorescens\_AA)} \org{P. fluorescens}\newline Price MN, Wetmore KM, Waters RJ, et al., Deutschbauer AM, Arkin AP. Nature 2018;557:503-509.\newline \href{https://doi.org/10.1038/s41586-018-0124-0}{doi:\allowbreak 10.\allowbreak 1038/\allowbreak s41586-\allowbreak 018-\allowbreak 0124-\allowbreak 0} & Cross-host fitness: part of the four-isolate within-species block, and the row that taught the most unannotated genes. & feba.\allowbreak db, org\allowbreak Id pseud\allowbreak o5\_\allowbreak N2C3\_\allowbreak 1 \\
\addlinespace[5pt]
57 & \textbf{Fitness Browser FW300-N1B4 (Aquipseudomonas sp.)} \org{P. fluorescens}\newline Price MN, Wetmore KM, Waters RJ, et al., Deutschbauer AM, Arkin AP. Nature 2018;557:503-509.\newline \href{https://doi.org/10.1038/s41586-018-0124-0}{doi:\allowbreak 10.\allowbreak 1038/\allowbreak s41586-\allowbreak 018-\allowbreak 0124-\allowbreak 0} & Cross-host fitness: completes the four-isolate within-species block, and names the case of a host with no species assignment. & feba.\allowbreak db, org\allowbreak Id pseud\allowbreak o1\_\allowbreak N1B4 \\
\addlinespace[5pt]
58 & \textbf{Fitness Browser B728a (Aquipseudomonas fragariae pv. syringae)} \org{P. syringae}\newline Price MN, Wetmore KM, Waters RJ, et al., Deutschbauer AM, Arkin AP. Nature 2018;557:503-509.\newline \href{https://doi.org/10.1038/s41586-018-0124-0}{doi:\allowbreak 10.\allowbreak 1038/\allowbreak s41586-\allowbreak 018-\allowbreak 0124-\allowbreak 0} & Cross-host fitness: the pathogen member of the block, and half of the one strain in it screened in two genetic backgrounds. & feba.\allowbreak db, org\allowbreak Id Syrin\allowbreak gaeB728a\allowbreak ; the GT\allowbreak DB renam\allowbreak e to A.\allowbreak  fragar\allowbreak iae pv.\allowbreak  syring\allowbreak ae means\allowbreak  both bi\allowbreak nomials \allowbreak are need\allowbreak ed as al\allowbreak iases \\
\addlinespace[5pt]
59 & \textbf{Fitness Browser B728a mexB deletion background} \org{P. syringae}\newline Price MN, Wetmore KM, Waters RJ, et al., Deutschbauer AM, Arkin AP. Nature 2018;557:503-509.\newline \href{https://doi.org/10.1038/s41586-018-0124-0}{doi:\allowbreak 10.\allowbreak 1038/\allowbreak s41586-\allowbreak 018-\allowbreak 0124-\allowbreak 0} & Cross-host fitness: the background-crossed half of the B728a pair, which is the only interaction design in the block. & feba.\allowbreak db, org\allowbreak Id Syrin\allowbreak gaeB728a\_\allowbreak mexBdel\allowbreak ta \\
\addlinespace[5pt]
60 & \textbf{Fitness Browser PA14 (Pseudomonas aeruginosa)} \org{P. aeruginosa}\newline Price MN, Wetmore KM, Waters RJ, et al., Deutschbauer AM, Arkin AP. Nature 2018;557:503-509.\newline \href{https://doi.org/10.1038/s41586-018-0124-0}{doi:\allowbreak 10.\allowbreak 1038/\allowbreak s41586-\allowbreak 018-\allowbreak 0124-\allowbreak 0} & Cross-host fitness: the same artifact and schema, kept visible because its size is the finding about this species rather than an omission. & feba.\allowbreak db, org\allowbreak Id PA14;\allowbreak  the bro\allowbreak wser doe\allowbreak s not re\allowbreak name thi\allowbreak s organi\allowbreak sm \\
\addlinespace[5pt]
\end{longtable}
\endgroup
\end{landscape}
